\section{Challenges and Outlook}
\label{sec:challenges}

\subsection{Open technical challenges}

\emph{Feedstock economics and supply.} Powder costs and critical-element exposure (Co, Ni, Ta, Nb) remain the largest barrier to marine adoption; WAAM wire routes mitigate cost but introduce evaporation and inclusion control problems \cite{cagirici2024,andersson2026waam}. Development of low-Co or Co-free corrosion-resistant HEAs is strategically important.

\emph{Defect--corrosion linkage.} The field lacks quantitative defect-tolerance criteria: how much porosity, of what size and connectivity, can a seawater-wetted AM HEA surface tolerate before pitting statistics become unacceptable? Such criteria exist implicitly for aerospace fatigue but not for marine corrosion, and they will be needed for classification \cite{dnvstb203}.

\emph{Environment-specificity of passivity.} The 316L-comparison evidence shows HEA passivity can collapse outside neutral chloride media \cite{lpbf_cantor_2024}. Crevice chemistry, under-biofilm acidification, and sulfides in stagnant zones all push toward those regimes and are untested for AM HEAs \cite{qiu2017}.

\emph{Coupled degradation modes.} Corrosion fatigue, SCC, cavitation, and erosion-corrosion data for AM HEAs are nearly absent, yet these govern marine component life \cite{lpbf_316l_scc2022}.

\emph{Biofouling--corrosion interaction.} Whether HEA surfaces, with their distinct passive-film chemistry, resist, tolerate, or accelerate biofilm-mediated attack is simply unknown; this connects the present topic to the companion literature on sustainable biofouling control \cite{nass2025marine}.

\subsection{Research directions}

\emph{Corrosion-by-design.} The demonstrated loop of thermodynamic design plus electrochemical validation \cite{lu2018design,quiambao2019} should be closed with machine-learning surrogate models trained on AM-specific microstructures, targeting compositions that are simultaneously printable (hot-crack resistant) and passive in seawater.

\emph{Process--passivity maps.} Systematic studies linking LPBF/WAAM parameter windows to passive-film composition (e.g., Cr/Mo enrichment measured by surface analysis) would convert corrosion from a screening outcome into a process-controlled property \cite{slm_ehea_2025,lpbf_coocrw2023}.

\emph{Real-seawater exposure networks.} Long-term immersion programs---tidal, splash, and submerged zones, with and without biofilm---for a small matrix of benchmark AM HEAs versus incumbent alloys, following the natural-exposure template already demonstrated for AM aluminum \cite{am_alsi10mg_2026}.

\emph{Hybrid structures.} Functionally graded DED builds and HEA claddings on marine steel, with galvanic and CP compatibility designed in, are the most economical near-term pathway \cite{jiang2019clad,cagirici2024}.

\emph{Digital qualification.} In-situ melt-pool monitoring plus build-data archives could support the data-driven material qualification frameworks classification societies are moving toward \cite{dnvstb203,abs_mam}, shortening the path for each new HEA composition.

\subsection{Deployment roadmap}

A realistic staged roadmap is: (1)~\emph{now}---benchmarking studies and repair cladding pilots on non-class-critical hardware; (2)~\emph{near term (2--5 years)}---qualified LPBF small components in seawater systems, WAAM propeller and bracket demonstrations, and the first natural-seawater exposure datasets; (3)~\emph{medium term (5--10 years)}---class-approved HEA compositions within marine AM standards, functionally graded marine structures, and life-cycle-cost validated substitution of incumbent alloys. Sustainability assessment should run alongside: HEA-AM parts must justify their critical-element burden through longer life, less maintenance, and recyclability to be credible in an industry moving toward life-cycle accountability \cite{nass2025marine}.
